% Image-free report with a shared BibTeX bibliography in references.bib.
% Build from reports/: pdflatex -synctex=1 -interaction=nonstopmode -halt-on-error main_report.tex
% Then run: bibtex main_report; pdflatex -synctex=1 main_report.tex (twice).
% This report contains every derivation; main_beamer.tex states the results.
\documentclass[11pt,a4paper]{article}
\usepackage[T1]{fontenc}
\usepackage[utf8]{inputenc}
\usepackage{lmodern,microtype}
\usepackage[margin=25mm]{geometry}
\usepackage{amsmath,amssymb,amsthm,mathtools}
\usepackage{booktabs,tabularx,array}
\usepackage[numbers,sort&compress]{natbib}
\usepackage[hidelinks]{hyperref}
\providecommand{\doi}[1]{\href{https://doi.org/#1}{\nolinkurl{doi:#1}}}
\setlength{\parskip}{4pt}
\setlength{\parindent}{0pt}
\setlength{\emergencystretch}{2em}
\newcommand{\ket}[1]{\lvert #1\rangle}
\newcommand{\bra}[1]{\langle #1\rvert}
\newcommand{\Tr}{\operatorname{Tr}}
\newcommand{\id}{\mathrm{id}}
\newcommand{\diag}{\operatorname{diag}}
\newcommand{\rank}{\operatorname{rank}}
\newcommand{\spec}{\operatorname{spec}}
\newcommand{\HH}{\mathcal H}
\newcommand{\ZZ}{\mathrm{ZZ}}
\newcommand{\ex}{\mathrm{ex}}
\newcommand{\OmF}{\Omega_{T_F}}
\newcommand{\OmR}{\Omega_R}
\newcommand{\dmat}{\mathsf d}
\newcommand{\Heff}{H_{\mathrm{eff}}}
\newtheorem{proposition}{Proposition}[section]
\newtheorem{lemma}[proposition]{Lemma}
\title{Two spins coupled to an XXZ bulk\\[3pt]
\large Reduced dynamics and quantum-state texture}
\author{Spin transport project}
\date{\today}

\begin{document}
\maketitle
\begin{abstract}
We study the joint reduced dynamics of two spin-$1/2$ probes,
each attached to one end of a finite XXZ bulk. The total Hamiltonian generates
unitary evolution, and the initial probe--bulk state is a product. For Ising
boundary coupling, four conditional bulk Hamiltonians give an exact
time-dependent Schur, or generalized dephasing, channel. Its complete action
is encoded in the output from the single probe preparation $\ket{++}$.
At every time where it is invertible, this channel has a canonical
time-local master equation with at most three nonzero decoherence rates
and diagonal Lindblad operators, all given in closed form by the
coherence factors of the Schur multiplier.
For exchange coupling, the available description gives the symmetry
blocks rather than closed expressions for their coefficients. When the Hamiltonian conserves total
magnetization and the bulk preparation commutes with bulk magnetization,
the reduced superoperator and Choi matrix split into blocks of sizes
$1,4,6,4,1$. Quantum-state texture is evaluated only on the output from
$\ket{++}$; exact relations to the Ising channel, its fidelities and its
entanglement-generating ability follow from the Schur structure.
\end{abstract}
\tableofcontents

\section{Introduction}

Two probes coupled to a common spin chain form an open quantum system
whose reduced dynamics depends on the interaction at each contact.
We consider Ising and exchange contacts separately, with no direct
probe--probe interaction. The bulk starts in a fixed state, uncorrelated
with the probe pair. Tracing out the bulk after unitary evolution defines
a quantum channel on the pair.

Ising contacts conserve both probe $z$ projections and produce generalized
dephasing in the joint computational basis. Exchange contacts allow spin
transfer between the probes and the bulk. Conservation of total
magnetization then constrains the reduced channel when the bulk
preparation respects the same rotation symmetry.

These two interactions lead to the following channel descriptions:
\begin{enumerate}
\item Ising coupling: an exact conditional-Hamiltonian reduction, an exact
Schur channel, and reconstruction of that channel from the full output
state for $\ket{++}$.
\item Ising coupling, time-local description: the canonical form of the
master equation, with its decoherence rates, Lindblad operators and
Hamiltonian expressed through the coherence factors.
\item Exchange coupling: exact total-magnetization sectors and, for a
magnetization-invariant bulk preparation, a phase-covariant channel with
five coherence-order blocks. These constraints leave the block entries
to be computed from the many-body evolution.
\item Texture: exact grand-sum, return-probability and response-operator
identities. For Ising coupling these yield channel fidelities and other
properties without additional input preparations.
\end{enumerate}

Section~\ref{sec:model} fixes the model, the reduced channel and the
texture conventions shared by both interactions. The two couplings are
then treated in separate parts: Part~\ref{part:ising} contains the Ising
channel (Section~\ref{sec:zz}) and its canonical form
(Section~\ref{sec:canonical}); Part~\ref{part:exchange} contains the
exchange channel (Section~\ref{sec:ex}).

The derivations use exact finite-system dynamics, without a weak-coupling
or Markov approximation. The channel structure fixes which probe
operators can mix; their time dependence depends on the bulk Hamiltonian
and its initial state.

\paragraph{Notation.}
To avoid collisions between symbols, we use the following conventions
throughout.
\begin{center}
\begin{tabular}{@{}ll@{}}
\toprule
Symbol & Meaning\\\midrule
$\Gamma(t)$, $\Gamma_{ab}(t)$ & Ising Schur matrix and its complex, dimensionless entries\\
$\eta_{ab}(t)=\dot\Gamma_{ab}/\Gamma_{ab}$ & logarithmic derivative of a coherence factor\\
$\kappa_{ab}$, $\omega_{ab}$ & coherence damping rate and phase rate, $\eta_{ab}=-\kappa_{ab}+i\omega_{ab}$\\
$\dmat=S+iA$ & $3\times3$ decoherence matrix on the diagonal operator sector\\
$\gamma_k(t)$ & canonical decoherence rates (eigenvalues of $\dmat$)\\
$F_\mu$, $G=(g^a_\mu)$ & diagonal traceless operators and their eigenvalue matrix\\
$\Heff=\diag(h_a)$ & canonical (effective) probe Hamiltonian\\
$\xi_{\pm2}$ & exchange-channel multiplier of the one-dimensional sectors $q=\pm2$\\
\bottomrule
\end{tabular}
\end{center}
Configuration labels $a,b,c,d$ run over $00,01,10,11$. The index
$\mu,\nu\in\{1,2,3\}$ labels $F_1,F_2,F_3$; $k$ labels canonical modes.
The couplings $g_A,g_B$ are unrelated to the eigenvalues $g^a_\mu$.

\section{Model, conventions and the reduced channel}
\label{sec:model}

\subsection{Geometry and Hamiltonians}

Let $N\ge2$ be the number of \emph{bulk} spins. There are $N+2$ spins in
total, with physical order $A,1,\ldots,N,B$. We use
\begin{equation}
 S^\alpha=\sigma^\alpha/2,\qquad
 S^\pm=S^x\pm iS^y=(\sigma^x\pm i\sigma^y)/2,\qquad \hbar=1.
\end{equation}
The basis convention is $\ket0=\ket\uparrow$, $\ket1=\ket\downarrow$,
and the probe basis is $(\ket{00},\ket{01},\ket{10},\ket{11})$.
For formulas involving a product of probe and bulk operators, we group
the Hilbert space as $\HH_{AB}\otimes\HH_E$ by a fixed permutation
of the physical site ordering.

The reference bulk and the two alternative interactions are
\begin{align}
 H_E&=J\sum_{j=1}^{N-1}
 (S_j^xS_{j+1}^x+S_j^yS_{j+1}^y+\Delta S_j^zS_{j+1}^z),
 \label{eq:bulk}\\
 V_{\ZZ}&=g_A S_A^zS_1^z+g_B S_N^zS_B^z,
 \label{eq:ising}\\
 V_{\ex}&=g_A(S_A^xS_1^x+S_A^yS_1^y)
          +g_B(S_N^xS_B^x+S_N^yS_B^y)\notag\\
 &=\frac{g_A}{2}(S_A^+S_1^-+S_A^-S_1^+)
   +\frac{g_B}{2}(S_N^+S_B^-+S_N^-S_B^+).
 \label{eq:exchange}
\end{align}
The second line of Eq.~\eqref{eq:exchange} follows from
$S^+S'^-+S^-S'^+=2(S^xS'^x+S^yS'^y)$ for commuting spins $S$ and $S'$.
Thus $H^{(r)}=H_E+V_r$, with $r=\ZZ$ or $\ex$.
Identity operators on the remaining spins are implicit. The parameters
$J,g_A,g_B$ have energy units and $\Delta$ is dimensionless. Equal
couplings mean $g_A=g_B=g$. The exchange flip-flop matrix element at
the left contact is $g_A/2$ in this convention.

Equation~\eqref{eq:bulk} has no applied fields. Longitudinal bulk or probe
fields preserve the magnetization arguments below, with their phases
included in the conditional dynamics. Transverse \emph{bulk} fields also
preserve the Ising Schur reduction, but generally invalidate the exchange
magnetization reduction because they break conservation of total
$z$ magnetization.

\subsection{Product preparation and exact channel}

Fix any bulk density matrix $\omega_E$, independent of the probe input.
The initial state and reduced evolution are
\begin{align}
 \rho(0)&=\rho_{AB}(0)\otimes\omega_E,\\
 U_t&=e^{-iH^{(r)}t},\qquad
 \Lambda_t(X)=\Tr_E[U_t(X\otimes\omega_E)U_t^\dagger],
 \label{eq:channel}\\
 \rho_{AB}(t)&=\Lambda_t(\rho_{AB}(0)).
\end{align}
The product assumption separates the pair from the bulk; it imposes no
factorization between $A$ and $B$ for the definition of the channel.
Only the texture protocol fixes their input to $\ket{++}$.
The bulk state $\omega_E$ may be pure or mixed, thermal or nonthermal,
and need not commute with $H_E$ or with the bulk magnetization unless a
specific result states otherwise. Changing $\omega_E$ changes the channel.

To see complete positivity and trace preservation directly, write
$\omega_E=\sum_\nu p_\nu\ket{\phi_\nu}\bra{\phi_\nu}$ with $p_\nu\ge0$,
$\sum_\nu p_\nu=1$, and choose any orthonormal bulk output basis
$\{\ket r\}$. Inserting $I_E=\sum_r\ket r\bra r$ into the partial trace,
\begin{equation}
 \Lambda_t(X)=\sum_{r,\nu}p_\nu\,{}_E\bra rU_t\ket{\phi_\nu}_E\,X\,
 {}_E\bra{\phi_\nu}U_t^\dagger\ket r_E .
\end{equation}
Each ${}_E\bra rU_t\ket{\phi_\nu}_E$ is an operator on $\HH_{AB}$. Hence
\begin{equation}
 K_{r\nu}(t)=\sqrt{p_\nu}\,{}_E\bra r U_t\ket{\phi_\nu}_E,
 \quad \Lambda_t(X)=\sum_{r,\nu}K_{r\nu}XK_{r\nu}^\dagger,
 \label{eq:kraus}
\end{equation}
and
\begin{equation}
 \sum_{r,\nu}K_{r\nu}^\dagger K_{r\nu}
 =\sum_\nu p_\nu\,{}_E\bra{\phi_\nu}U_t^\dagger
 \Bigl(\sum_r\ket r\bra r\Bigr)U_t\ket{\phi_\nu}_E
 =\sum_\nu p_\nu\,\langle\phi_\nu|\phi_\nu\rangle\,I_4=I_4 .
 \label{eq:krausTP}
\end{equation}
These are formal exact expressions in both regimes; evaluating them for
an interacting bulk remains a many-body problem. Correlated probe--bulk
preparations have reduced trajectories, but do not in general define
the same CPTP map on arbitrary probe inputs.

\subsection{Choi normalization}

The \emph{normalized} Choi state, with the output factor first, is
\begin{equation}
 \ket{\Phi_4}=\frac12\sum_{a=1}^4\ket a_{AB}\ket a_R,
 \qquad J_t=(\Lambda_t\otimes\id_R)(\ket{\Phi_4}\bra{\Phi_4}).
 \label{eq:choi}
\end{equation}
Since $\ket{\Phi_4}\bra{\Phi_4}=\frac14\sum_{a,b}E_{ab}\otimes E_{ab}$
with $E_{ab}=\ket a\bra b$,
\begin{equation}
 J_t=\frac14\sum_{a,b}\Lambda_t(E_{ab})\otimes E_{ab}.
 \label{eq:choigeneral}
\end{equation}
It obeys $J_t\succeq0$, $\Tr J_t=1$ and
$\Tr_{AB}J_t=I_R/4$. The corresponding unnormalized Choi matrix is
$C_t=4J_t$. The auxiliary system $R$ represents a reference for the
channel and is independent of the physical bulk.

\subsection{Quantum-state texture of the \texorpdfstring{$\ket{++}$}{++} preparation}
\label{sec:texture}

Texture is defined relative to the probe computational basis. Its
textureless state, and the sole probe preparation used for texture, is
\begin{equation}
 \ket{f_1}=\ket+_A\otimes\ket+_B
 =\frac12\sum_{a=1}^4\ket a,\qquad
 f_1=\ket{f_1}\bra{f_1},\qquad (f_1)_{ab}=\frac14.
 \label{eq:f1}
\end{equation}
Denote its full output state by $\rho_f(t)=\Lambda_t(f_1)$. For any
probe state, the grand sum is real and satisfies
\begin{equation}
 \Sigma(\rho)=\sum_{a,b}\rho_{ab}=4\bra{f_1}\rho\ket{f_1},
 \qquad 0\le\Sigma(\rho)\le4.
\end{equation}
The equality follows from $\bra{f_1}\rho\ket{f_1}=\frac14\sum_{a,b}\rho_{ab}$;
the bounds follow because $\bra{f_1}\rho\ket{f_1}$ is a probability.
For the single physical input $f_1\otimes\omega_E$, set
\begin{equation}
 p_f(t)=\bra{f_1}\rho_f(t)\ket{f_1}=\Sigma(\rho_f(t))/4,
 \quad T_F(\rho_f)=1-p_f,\quad R(\rho_f)=-\ln p_f.
 \label{eq:textures}
\end{equation}
These are the fidelity texture and rugosity conventions used here
\cite{Parisio2024,Vieira2026}. At $p_f=0$, $R=+\infty$; initially
$p_f=1$ and both textures vanish. The texturizing powers evaluated on
the unique textureless input are $\OmF(t)=T_F[\Lambda_t(f_1)]$ and
$\OmR(t)=R[\Lambda_t(f_1)]$.

Since $\ket+\bra+=(I+\sigma^x)/2$, we have $4f_1=(I+\sigma_A^x)(I+\sigma_B^x)$, and
\begin{equation}
 \Sigma(\rho_f)=\Tr[4f_1\rho_f]=1+\langle\sigma_A^x\rangle
 +\langle\sigma_B^x\rangle+\langle\sigma_A^x\sigma_B^x\rangle.
 \label{eq:measurement}
\end{equation}
A joint $x$-basis measurement gives $p_f=P_{++}^{(x)}$ directly.
Equivalently, apply a Hadamard to each probe and record outcome $00$.
This uses one measurement setting, whereas full channel reconstruction
in the Ising case uses the additional coherence measurements described
in Section~\ref{sec:oneinput}.

For any channel, the dual map $\Lambda_t^*$, defined by
$\Tr[Y\Lambda_t(X)]=\Tr[\Lambda_t^*(Y)X]$, gives the grand-sum response operator
\begin{equation}
 \mathcal R_t=4\Lambda_t^*(f_1),\qquad
 \Sigma[\Lambda_t(X)]=\Tr[4f_1\Lambda_t(X)]=\Tr(\mathcal R_tX).
 \label{eq:response}
\end{equation}
These definitions apply in both coupling regimes. Their consequences
for each interaction are derived in Sections~\ref{sec:zztexture}
and~\ref{sec:extexture}.

\part{Ising coupling}
\label{part:ising}

\section{Ising coupling: exact generalized dephasing}
\label{sec:zz}

\subsection{Four conditional Hamiltonians}

The probe operators $S_A^z,S_B^z$ commute with $H_E$ (which acts only on
$E$) and with $V_{\ZZ}$ (which contains only $z$ operators on the probes).
Hence $[H^{(\ZZ)},S_A^z]=[H^{(\ZZ)},S_B^z]=0$. Label each probe basis
state by $a=(s_A^a,s_B^a)$ with $s_A^a,s_B^a=\pm1/2$, and set
$P_a=\ket a\bra a$. The projectors satisfy
\begin{equation}
 P_aP_b=\delta_{ab}P_a,\qquad \sum_aP_a=I_{AB},\qquad
 S_A^z=\sum_as_A^aP_a,\qquad S_B^z=\sum_as_B^aP_a .
\end{equation}
Inserting these resolutions into
$H^{(\ZZ)}=I_{AB}\otimes H_E+g_AS_A^z\otimes S_1^z+g_BS_B^z\otimes S_N^z$ gives
\begin{equation}
 H^{(\ZZ)}=\sum_a P_a\otimes H_a,\qquad
 H_a=H_E+g_A s_A^a S_1^z+g_B s_B^a S_N^z.
 \label{eq:conditional}
\end{equation}
By induction, $(H^{(\ZZ)})^n=\sum_aP_a\otimes H_a^n$: if it holds for $n$,
then $(H^{(\ZZ)})^{n+1}=\sum_{a,b}P_aP_b\otimes H_a^nH_b=\sum_aP_a\otimes H_a^{n+1}$.
Summing the exponential series term by term,
\begin{equation}
 U_t=\sum_aP_a\otimes U_a(t),\qquad U_a(t)=e^{-iH_at}.
 \label{eq:controlled}
\end{equation}
This is exact for any bulk Hamiltonian acting only on $E$, including one
with transverse fields.

For equal couplings, the four perturbations are
\begin{equation}
\begin{array}{c|c}
 a&H_a-H_E\\\hline
 00&\frac g2(S_1^z+S_N^z)\\
 01&\frac g2(S_1^z-S_N^z)\\
 10&-\frac g2(S_1^z-S_N^z)\\
 11&-\frac g2(S_1^z+S_N^z)
\end{array}.
\end{equation}
These configurations label invariant branches. The input may still
coherently occupy all four branches. In general, the differences
\begin{equation}
 H_a-H_b=g_A(s_A^a-s_A^b)S_1^z+g_B(s_B^a-s_B^b)S_N^z
 \label{eq:difference}
\end{equation}
involve only the two bulk boundary operators.

\subsection{Reduced state for an arbitrary probe input}
\label{sec:reducedstate}

Take an arbitrary probe state and the fixed bulk state,
\begin{equation}
 \rho_{AB}(0)=\sum_{a,b}r_{ab}E_{ab},\qquad
 \rho_{ABE}(0)=\sum_{a,b}r_{ab}\,E_{ab}\otimes\omega_E .
\end{equation}
Inserting Eq.~\eqref{eq:controlled} on both sides,
\begin{equation}
 \rho_{ABE}(t)=U_t\rho_{ABE}(0)U_t^\dagger
 =\sum_{c,d}\sum_{a,b}r_{ab}\,(P_cE_{ab}P_d)\otimes(U_c\omega_EU_d^\dagger).
\end{equation}
Since $P_cE_{ab}P_d=\delta_{ca}\delta_{db}E_{ab}$,
\begin{equation}
 \rho_{ABE}(t)=\sum_{a,b}r_{ab}\,E_{ab}\otimes U_a(t)\,\omega_E\,U_b(t)^\dagger .
 \label{eq:jointstate}
\end{equation}
The bulk operator multiplying $E_{ab}$ depends on the ordered pair
$(a,b)$, although the initial bulk state is the same for all pairs.
Taking the partial trace term by term,
\begin{equation}
 \rho_{AB}(t)=\sum_{a,b}r_{ab}\,\Gamma_{ab}(t)\,E_{ab},\qquad
 \Gamma_{ab}(t)=\Tr_E[U_a(t)\omega_EU_b(t)^\dagger]
 =\Tr_E[\omega_EU_b(t)^\dagger U_a(t)],
 \label{eq:gamma}
\end{equation}
where the last form uses cyclicity of the trace on $E$. In the
computational order $(00,01,10,11)$ the full output is
\begin{equation}
\rho_{AB}(t)=
\begin{pmatrix}
 r_{00,00}&r_{00,01}\Gamma_{00,01}&r_{00,10}\Gamma_{00,10}&r_{00,11}\Gamma_{00,11}\\
 r_{01,00}\Gamma_{01,00}&r_{01,01}&r_{01,10}\Gamma_{01,10}&r_{01,11}\Gamma_{01,11}\\
 r_{10,00}\Gamma_{10,00}&r_{10,01}\Gamma_{10,01}&r_{10,10}&r_{10,11}\Gamma_{10,11}\\
 r_{11,00}\Gamma_{11,00}&r_{11,01}\Gamma_{11,01}&r_{11,10}\Gamma_{11,10}&r_{11,11}
\end{pmatrix},
\label{eq:sixteen}
\end{equation}
where $\Gamma_{aa}=1$ and $\Gamma_{ba}=\Gamma_{ab}^*$ have been used
(Section~\ref{sec:schurconstraints}). All populations $r_{aa}$ are
unchanged; only the six independent complex coherences evolve.

For a pure bulk $\omega_E=\ket{\phi_0}\bra{\phi_0}$ and a pure probe input
$\sum_ac_a\ket a$, the global state is
$\sum_ac_a\ket a\otimes U_a\ket{\phi_0}$ and
$\Gamma_{ab}=\bra{\phi_0}U_b^\dagger U_a\ket{\phi_0}$ is the complex overlap
of two conditional bulk histories. Its squared modulus is a conditional
Loschmidt echo; the complex amplitude, including its phase, is needed
for the channel.

\subsection{Schur form and its physical constraints}
\label{sec:schurconstraints}

Since Eq.~\eqref{eq:gamma} holds for every input, the channel acts on
each matrix unit as
\begin{equation}
 \boxed{\Lambda_t(E_{ab})=\Gamma_{ab}(t)E_{ab},\qquad
 \Lambda_t(X)=\Gamma(t)\circ X,}
 \label{eq:schur}
\end{equation}
where $\circ$ is the entrywise Schur (Hadamard) product and
$\Gamma(t)=(\Gamma_{ab}(t))_{a,b=1}^4$. The sixteen $E_{ab}$ are
superoperator eigenoperators, with eigenvalues $\Gamma_{ab}$, four of
which equal one.

The matrix $\Gamma$ is a correlation matrix:
\begin{equation}
 \Gamma=\Gamma^\dagger\succeq0,\qquad \Gamma_{aa}=1,
 \qquad |\Gamma_{ab}|\le1,\qquad \Gamma_{ab}(0)=1.
 \label{eq:correlation}
\end{equation}
\emph{Proof.} Set $A_a(t)=U_a(t)\sqrt{\omega_E}$, so that
$\Gamma_{ab}=\Tr_E(A_aA_b^\dagger)$. For any $v\in\mathbb C^4$,
\begin{equation}
 v^\dagger\Gamma v=\sum_{a,b}v_a^*\Gamma_{ab}v_b
 =\Tr_E\Bigl[\Bigl(\sum_av_a^*A_a\Bigr)\Bigl(\sum_bv_b^*A_b\Bigr)^\dagger\Bigr]\ge0,
\end{equation}
because $(\sum_bv_b^*A_b)^\dagger=\sum_bv_bA_b^\dagger$ and $\Tr(YY^\dagger)\ge0$.
Hence $\Gamma=\Gamma^\dagger\succeq0$. Unitarity and $\Tr\omega_E=1$
give $\Gamma_{aa}=\Tr_E(U_a\omega_EU_a^\dagger)=1$. Nonnegativity of each
$2\times2$ principal minor gives
$\Gamma_{aa}\Gamma_{bb}-|\Gamma_{ab}|^2=1-|\Gamma_{ab}|^2\ge0$.
Finally $U_a(0)=I$ gives $\Gamma_{ab}(0)=1$. Equivalently, for a
decomposition $\omega_E=\sum_\nu p_\nu\ket{\phi_\nu}\bra{\phi_\nu}$ and an
output basis $\{\ket r\}$, the numbers
$k_{r\nu,a}=\sqrt{p_\nu}\bra rU_a\ket{\phi_\nu}$ give the Gram
factorization $\Gamma_{ab}=\sum_{r,\nu}k_{r\nu,a}k_{r\nu,b}^*$.\qed

Conversely, every positive matrix with unit diagonal defines a CPTP
Schur map, by the Kraus construction of Eq.~\eqref{eq:diagkraus}. This
converse characterizes the map class, not which members a specified XXZ
Hamiltonian realizes at a given time.

All four populations are fixed and $\Lambda_t(I_4)=\sum_a\Gamma_{aa}E_{aa}=I_4$,
so the channel is unital. Six complex off-diagonal factors specify the
nontrivial action, subject to positivity. Because the superoperator is
diagonal in the Hilbert--Schmidt orthonormal basis $\{E_{ab}\}$, its
singular values are $|\Gamma_{ab}|$. The common interacting bulk can
correlate the two probes, so a product of independent single-probe
dephasing channels cannot be assumed. There is no general monotonicity
of $|\Gamma_{ab}(t)|$.

\subsection{Exact expressions for the coherence factors and their derivatives}
\label{sec:gammaderivative}

For any convex decomposition $\omega_E=\sum_\nu p_\nu\ket{\phi_\nu}\bra{\phi_\nu}$,
Eq.~\eqref{eq:gamma} reads
\begin{equation}
 \Gamma_{ab}(t)=\sum_\nu p_\nu
 \bra{\phi_\nu}U_b(t)^\dagger U_a(t)\ket{\phi_\nu}.
 \label{eq:gammaensemble}
\end{equation}
With $\dot U_a=-iH_aU_a$ and $\dot U_b^\dagger=iU_b^\dagger H_b$
(both using $[H_a,U_a]=[H_b,U_b]=0$),
\begin{equation}
 \dot\Gamma_{ab}
 =\Tr_E\bigl[\omega_E(iU_b^\dagger H_b)U_a+\omega_EU_b^\dagger(-iH_aU_a)\bigr]
 =-i\Tr_E\bigl[\omega_EU_b(t)^\dagger(H_a-H_b)U_a(t)\bigr].
 \label{eq:gammadot}
\end{equation}
With Eq.~\eqref{eq:difference},
\begin{equation}
 \dot\Gamma_{ab}=-i\Bigl[g_A(s_A^a-s_A^b)\Tr_E(\omega_EU_b^\dagger S_1^zU_a)
 +g_B(s_B^a-s_B^b)\Tr_E(\omega_EU_b^\dagger S_N^zU_a)\Bigr].
 \label{eq:gammadotboundary}
\end{equation}
Only the two boundary operators enter. No numerical time
differentiation is required once the conditional states and the matrix
elements of $S_1^z$ and $S_N^z$ are available.

At $t=0$, $U_a(0)=I$ gives
\begin{equation}
 \Gamma_{ab}(0)=1,\qquad
 \dot\Gamma_{ab}(0)=-i(\varepsilon_a-\varepsilon_b),\qquad
 \varepsilon_a=\Tr_E(\omega_EH_a)\in\mathbb R .
 \label{eq:initialderivative}
\end{equation}
The initial derivative is purely imaginary. Its consequence for the
canonical generator is derived in Section~\ref{sec:initial}.

\subsection{Dephasing and depolarization}\label{sec:depolarizing}

A depolarizing map has the form
$\mathcal D_\lambda(X)=\lambda X+(1-\lambda)\Tr(X)I_4/4$.
The Ising map fixes every $P_a$. Requiring
$\mathcal D_\lambda(P_a)=\lambda P_a+(1-\lambda)I_4/4=P_a$ gives
$(1-\lambda)(I_4/4-P_a)=0$, which forces $\lambda=1$.
Thus a nonidentity Ising Schur map cannot be a depolarizing map.
Dephasing includes coherent diagonal rotations and revivals; it does
not imply irreversible or strictly dissipative evolution.

\subsection{One output determines the complete channel}
\label{sec:oneinput}

For the texture preparation $f_1$ of Eq.~\eqref{eq:f1}, with full
output state $\rho_f(t)=\Lambda_t(f_1)$, the Schur form gives
$[\rho_f(t)]_{ab}=\Gamma_{ab}(t)(f_1)_{ab}=\Gamma_{ab}(t)/4$, hence
\begin{equation}
 \boxed{\rho_f(t)=\Gamma(t)/4,\qquad
 \Gamma(t)=4\rho_f(t),\qquad
 \Lambda_t(X)=[4\rho_f(t)]\circ X.}
 \label{eq:reconstruction}
\end{equation}
Consequently the \emph{entire density matrix} of one known output
determines every property of the instantaneous Ising channel. Temporal
properties require these outputs at the relevant times. A single scalar
texture value cannot replace this matrix. This identity is special to
the Ising Schur family; it does not hold for arbitrary two-probe channels.

The populations of $\rho_f$ are $1/4$, so its $12$ remaining real
parameters are its six complex coherences. In a Pauli expansion,
$\langle ZI\rangle=\langle IZ\rangle=\langle ZZ\rangle=0$.
The eight local measurement settings $(\alpha,\beta)\ne(z,z)$,
$\alpha,\beta\in\{x,y,z\}$, recover the other twelve Pauli coefficients,
including single-spin expectations from their marginals. Thus one
input preparation together with output tomography suffices.

The time-local data of Section~\ref{sec:canonical} follow from the same
output and its time derivative:
\begin{equation}
 \dot\Gamma(t)=4\dot\rho_f(t),\qquad
 \eta_{ab}(t)=\frac{\dot\Gamma_{ab}(t)}{\Gamma_{ab}(t)}
 =\frac{[\dot\rho_f(t)]_{ab}}{[\rho_f(t)]_{ab}}
 \quad(a\ne b,\ [\rho_f(t)]_{ab}\ne0).
 \label{eq:etafromrho}
\end{equation}

\subsection{Choi, Kraus and spectral consequences}
\label{sec:choikraus}

Inserting Eq.~\eqref{eq:schur} into Eq.~\eqref{eq:choigeneral},
\begin{equation}
 J_t=\frac14\sum_{a,b}\Gamma_{ab}\,E_{ab}\otimes E_{ab}
 =\frac14\sum_{a,b}\Gamma_{ab}\ket{aa}\bra{bb}.
\end{equation}
Let $W:\HH_{AB}\to\HH_{AB}\otimes\HH_R$ be the isometry $W\ket a=\ket{aa}$.
Then $W(\Gamma/4)W^\dagger=\frac14\sum_{a,b}\Gamma_{ab}\ket{aa}\bra{bb}$, so
\begin{equation}
 J_t=W\frac{\Gamma(t)}4W^\dagger=W\rho_f(t)W^\dagger.
 \label{eq:choizz}
\end{equation}
Since $W^\dagger W=I_4$, conjugation by $W$ preserves nonzero spectra.
The Choi support has dimension at most four. In particular,
\begin{align}
 \rank J_t&=\rank\Gamma\le4,\\
 \spec(J_t)&=\spec(\Gamma/4)\ \text{with twelve additional zeros},\\
 S(J_t)&=S(\rho_f),\qquad
 \Tr J_t^2=\Tr\rho_f^2=\frac1{16}\sum_{a,b}|\Gamma_{ab}|^2.
 \label{eq:choipurity}
\end{align}

If $\Gamma=\sum_k\lambda_kv_kv_k^\dagger$ with $\lambda_k\ge0$ and
normalized eigenvectors $v_k$, one minimal Kraus family is
\begin{equation}
 D_k=\sqrt{\lambda_k}\,\diag(v_{k,1},\ldots,v_{k,4}),
 \qquad \Lambda_t(X)=\sum_k D_kXD_k^\dagger.
 \label{eq:diagkraus}
\end{equation}
Indeed $\sum_kD_kE_{ab}D_k^\dagger=\sum_k\lambda_kv_{k,a}v_{k,b}^*E_{ab}
=\Gamma_{ab}E_{ab}$, and $\sum_kD_k^\dagger D_k=\diag(\Gamma_{aa})=I_4$.
Choi rank and superoperator rank are distinct: the identity has Choi
rank one and superoperator rank sixteen.

The Kraus decomposition diagonalizes the positive $4\times4$ matrix
$\Gamma$. The canonical time-local generator of Section~\ref{sec:canonical}
requires a \emph{different} diagonalization, that of a $3\times3$
decoherence matrix built from logarithmic derivatives of $\Gamma$.
The eigenvalues of $\Gamma$ are not the canonical decoherence rates.

\paragraph{Unitarity.}
The map is unitary iff $\rank\Gamma=1$. If $\rank\Gamma=1$, then
$\Gamma=\lambda vv^\dagger$; unit diagonal gives $\lambda|v_a|^2=1$, so
$u_a=\sqrt\lambda\,v_a$ obeys $|u_a|=1$ and $\Gamma_{ab}=u_au_b^*$. Then
$\Lambda_t(X)=\diag(u)X\diag(u)^\dagger$. Conversely, a unitary channel has
a pure Choi state, and by Eq.~\eqref{eq:choizz} $\rank\Gamma=\rank J_t=1$.
This does not require that the unitary factor into two local rotations.

\paragraph{Reference entanglement.}
Take the partial transpose of $J_t$ on $R$:
$\ket{aa}\bra{bb}\mapsto\ket{ab}\bra{ba}$, so
$J_t^{T_R}=\frac14\sum_{a,b}\Gamma_{ab}\ket{ab}\bra{ba}$. The terms with
$a=b$ give the eigenvalue $1/4$ four times, on $\ket{aa}$. For each pair
$a<b$, the two-dimensional block on $\{\ket{ab},\ket{ba}\}$ is
\begin{equation}
 \frac14\begin{pmatrix}0&\Gamma_{ab}\\\Gamma_{ab}^*&0\end{pmatrix},
 \qquad\text{eigenvalues }\pm\frac{|\Gamma_{ab}|}{4}.
\end{equation}
Hence its reference--output negativity is
\begin{equation}
 \mathcal N_{AB|R}(J_t)=\frac14\sum_{a<b}|\Gamma_{ab}|.
\end{equation}
This is distinct from physical $A|B$ entanglement. The Ising map is
entanglement breaking iff all off-diagonal $\Gamma_{ab}$ vanish:
otherwise its Choi state has a negative partial transpose and is
entangled, while the diagonal case has the explicitly separable Choi
state $\frac14\sum_a\ket{aa}\bra{aa}$.

\subsection{Invertibility, time-local generator and intermediate maps}
\label{sec:invertibility}

Because the superoperator is diagonal on $\{E_{ab}\}$,
\begin{equation}
 \det\Lambda_t=\prod_{a,b}\Gamma_{ab}(t)
 =\prod_{a<b}\Gamma_{ab}\Gamma_{ba}=\prod_{a<b}|\Gamma_{ab}(t)|^2 ,
\end{equation}
using $\Gamma_{aa}=1$ and $\Gamma_{ba}=\Gamma_{ab}^*$. Therefore
\begin{equation}
 \Lambda_t\ \text{invertible}\iff\Gamma_{ab}(t)\ne0\quad\forall a<b,
 \qquad \Lambda_t^{-1}(E_{ab})=\Gamma_{ab}(t)^{-1}E_{ab}.
\end{equation}
At such times the exact time-local generator
$\mathcal L_t=\dot\Lambda_t\Lambda_t^{-1}$ acts as
\begin{equation}
 \mathcal L_t(E_{ab})=\frac{\dot\Gamma_{ab}(t)}{\Gamma_{ab}(t)}E_{ab}
 \equiv\eta_{ab}(t)E_{ab}.
 \label{eq:generatorschur}
\end{equation}

If $\Lambda_s$ is invertible, the unique intermediate linear map
$\Lambda_{t,s}$ with $\Lambda_t=\Lambda_{t,s}\Lambda_s$ is
$\Lambda_{t,s}=\Lambda_t\Lambda_s^{-1}$, i.e.\ Schur multiplication by
$Q_{ab}(t,s)=\Gamma_{ab}(t)/\Gamma_{ab}(s)$. The matrix $Q$ is Hermitian,
since $Q_{ba}=\Gamma_{ab}^*(t)/\Gamma_{ab}^*(s)=Q_{ab}^*$, and has unit
diagonal. A Schur map with unit diagonal is trace preserving, and it is
completely positive iff its multiplier is positive semidefinite (its
normalized Choi state is $WQW^\dagger/4$). Hence
\begin{equation}
 \Lambda_{t,s}\text{ is CPTP}\quad\Longleftrightarrow\quad
 Q(t,s)\succeq0.
 \label{eq:divisibility}
\end{equation}
At singular times the division is undefined and this test cannot be
applied.

\begin{lemma}\label{lem:nosingular}
On an interval where the coherence factors obey
$\dot\Gamma_{ab}=\eta_{ab}\Gamma_{ab}$ with locally integrable $\eta_{ab}$,
no $\Gamma_{ab}$ can vanish at a finite time if it is nonzero at the start
of the interval.
\end{lemma}
\begin{proof}
The linear equation integrates to
$\Gamma_{ab}(t)=\Gamma_{ab}(s)\exp\int_s^t\eta_{ab}(\tau)\,d\tau$, and the
exponential of a finite integral is nonzero. Equivalently,
$\det\Lambda_t=\det\Lambda_s\exp\int_s^t\Tr\mathcal L_\tau\,d\tau$.
\end{proof}
Thus a finite, locally integrable time-local generator cannot take an
invertible evolution to a singular channel in finite time. If a
coherence factor vanishes at $s$ and later becomes nonzero, no linear map
$\Lambda_{t,s}$ on the probes can satisfy $\Lambda_t=\Lambda_{t,s}\Lambda_s$,
because $\Lambda_s(E_{ab})=0$ while $\Lambda_t(E_{ab})\ne0$. This
precludes divisibility through that time. The global unitary and the
reduced CPTP maps $\Lambda_t$ themselves remain analytic in $t$.

\subsection{Probe entanglement from the same input}
\label{sec:entangling}

For two-qubit states, concurrence $C$ and negativity $\mathcal N$ test
entanglement of $A$ with $B$; their definitions are in
Appendix~\ref{app:observables}. For a Schur channel, the output from
$\ket{++}$ also determines the maximum concurrence generated from
separable probe states.

\begin{lemma}[Local-filter identity]\label{lem:filter}
For $2\times2$ matrices $A,B$ and a positive two-qubit operator $\rho$,
with concurrence defined from the eigenvalues of $\rho\widetilde\rho$
without normalization,
$C[(A\otimes B)\rho(A\otimes B)^\dagger]=|\det A\det B|\,C(\rho)$.
\end{lemma}
\begin{proof}
For any $2\times2$ matrix, $A^{\mathsf T}\sigma_yA=\det(A)\,\sigma_y$.
Taking the complex conjugate and using $\sigma_y^*=-\sigma_y$ gives
$A^\dagger\sigma_yA^*=\det(A)^*\sigma_y$. For invertible $A$, the first
identity also gives $\sigma_yA^{\mathsf T}\sigma_y=\det(A)A^{-1}$. Let
$\rho'=(A\otimes B)\rho(A\otimes B)^\dagger$ and $Y=\sigma_y\otimes\sigma_y$.
Then
\begin{align}
 \rho'\widetilde{\rho'}
 &=(A\otimes B)\rho(A\otimes B)^\dagger Y(A\otimes B)^*\rho^*(A\otimes B)^{\mathsf T}Y\notag\\
 &=\det(A)^*\det(B)^*\,(A\otimes B)\rho\,Y\rho^*Y\,Y(A\otimes B)^{\mathsf T}Y\notag\\
 &=|\det A\det B|^2\,(A\otimes B)\,\rho\widetilde\rho\,(A\otimes B)^{-1}.
\end{align}
This is similar to $|\det A\det B|^2\rho\widetilde\rho$, so the square
roots of its eigenvalues, and hence $C$, scale by $|\det A\det B|$.
Singular filters follow by continuity of $C$.
\end{proof}

\begin{proposition}
For a fixed Ising Schur channel, its maximum output concurrence over
separable probe inputs is
\begin{equation}
 \max_{\rho\in\mathrm{Sep}(A:B)}C[\Lambda_t(\rho)]
 =C[\rho_f(t)]=C[\Gamma(t)/4].
 \label{eq:cgen}
\end{equation}
\end{proposition}
\begin{proof}
For a normalized pure product input
$\ket\psi=(a_0\ket0+a_1\ket1)\otimes(b_0\ket0+b_1\ket1)$, the components
are $\psi_{ij}=a_ib_j$. By Eq.~\eqref{eq:schur},
$\Lambda_t(\ket\psi\bra\psi)_{ab}=\psi_a\Gamma_{ab}\psi_b^*$, i.e.
$\Lambda_t(\ket\psi\bra\psi)=D_\psi\Gamma D_\psi^\dagger$ with
$D_\psi=\diag(\psi)$. Using $\Gamma=4\rho_f$ and $2D_\psi=A\otimes B$ with
$A=\sqrt2\diag(a_0,a_1)$, $B=\sqrt2\diag(b_0,b_1)$,
\begin{equation}
 \Lambda_t(\ket\psi\bra\psi)=(A\otimes B)\rho_f(A\otimes B)^\dagger .
\end{equation}
This output has unit trace because $\rho_f$ has uniform diagonal $1/4$:
its trace is $\sum_{ij}4|a_i|^2|b_j|^2/4=1$. Lemma~\ref{lem:filter} gives
$C_{\mathrm{out}}=|\det A\det B|C(\rho_f)=4|a_0a_1b_0b_1|C(\rho_f)\le C(\rho_f)$,
since $|a_0a_1|\le\frac12(|a_0|^2+|a_1|^2)=\frac12$ and likewise for $b$.
A separable input is a convex combination
$\rho=\sum_iq_i\ket{\psi_i}\bra{\psi_i}$ of pure product states.
Linearity of $\Lambda_t$ and convexity of $C$ give
$C[\Lambda_t(\rho)]\le\sum_iq_iC[\Lambda_t(\ket{\psi_i}\bra{\psi_i})]\le C(\rho_f)$.
The input $\ket{++}$ attains equality.
\end{proof}
In particular, the Ising map can create probe entanglement from some
separable input iff $\rho_f(t)$ is entangled. Texture and Choi
reference--output negativity are different quantities and do not replace
this criterion.

\subsection{Ising texture and channel fidelities}
\label{sec:zztexture}

Using $\rho_f=\Gamma/4$ in Eq.~\eqref{eq:textures}, together with
$\sum_{a,b}\Gamma_{ab}=4+\sum_{a<b}(\Gamma_{ab}+\Gamma_{ab}^*)
=4+2\sum_{a<b}\Re\Gamma_{ab}$,
\begin{align}
 p_f&=\frac1{16}\sum_{a,b}\Gamma_{ab}
     =\frac14+\frac18\sum_{a<b}\Re\Gamma_{ab},\\
 \OmF&=\frac34-\frac18\sum_{a<b}\Re\Gamma_{ab},
 &\OmR&=-\ln\!\left(\frac1{16}\sum_{a,b}\Gamma_{ab}\right).
 \label{eq:omegazz}
\end{align}
For the same Schur channel, the entanglement fidelity relative to the
identity is
\begin{equation}
 F_e=\bra{\Phi_4}J_t\ket{\Phi_4}
 =\frac14\sum_{c,d}\bra{cc}\Bigl(\frac14\sum_{a,b}\Gamma_{ab}\ket{aa}\bra{bb}\Bigr)\ket{dd}
 =\frac1{16}\sum_{a,b}\Gamma_{ab}=p_f .
\end{equation}
The general relation $F_{\mathrm{avg}}=(dF_e+1)/(d+1)$ between average
and entanglement fidelities \cite{Nielsen2002}, with $d=4$, then gives
\begin{equation}
 \boxed{F_e=p_f=\Sigma(\rho_f)/4,\qquad
 F_{\mathrm{avg}}=\frac{4F_e+1}{5}
                 =\frac{1+\Sigma(\rho_f)}5.}
 \label{eq:fidelities}
\end{equation}
The equality $F_e=p_f$ uses the Schur structure.
These are fidelities of the \emph{two-probe map relative to the identity},
not fidelities of sending an unknown state from $A$ to $B$.

Only the real sum of the six coherences enters $p_f$. It does not fix
their individual magnitudes or phases and does not determine purity,
probe entanglement, or CP divisibility. Coherent phase evolution can
change texture even when the map is unitary.

\subsection{Response operator and texture preservation}

In the Ising case, Eq.~\eqref{eq:response} reads
$\Sigma[\Lambda_t(X)]=\sum_{a,b}\Gamma_{ab}X_{ab}$, while
$\Tr(\mathcal R_tX)=\sum_{a,b}(\mathcal R_t)_{ba}X_{ab}$. Comparing
coefficients of $X_{ab}$ gives $(\mathcal R_t)_{ba}=\Gamma_{ab}$, i.e.
\begin{equation}
 \mathcal R_t=\Gamma(t)^{\mathsf T}.
 \label{eq:responsetranspose}
\end{equation}
The transpose is ordinary, not a Hermitian conjugate; since $\Gamma$ is
Hermitian, $\Gamma^{\mathsf T}=\Gamma^*$. In general, fixing the free
state, $\Lambda_t(f_1)=f_1$, and preserving the grand sum for every input,
$\Lambda_t^*(f_1)=f_1$, are different conditions \cite{Vieira2026}. For
this Schur family, $\Lambda_t(f_1)=\Gamma/4$ and
$\Lambda_t^*(f_1)=\Gamma^{\mathsf T}/4$, so each condition forces
$\Gamma_{ab}=1$ for all $a,b$, hence the identity channel.
Moreover $\OmF=0$ means $p_f=1$, i.e.\ $\sum_{a<b}\Re\Gamma_{ab}=6$; with
$|\Gamma_{ab}|\le1$ this forces every $\Gamma_{ab}=1$. Consequently
$\OmF=0$ also certifies the identity within the known Ising family.

\subsection{Summary of the Ising channel structure}

Six complex overlaps between conditional bulk evolutions determine the
Ising channel. The output $\rho_f=\Gamma/4$ therefore specifies its full
action, including its Choi state, fidelities and ability to generate
probe entanglement. Evaluating the overlaps requires the conditional
dynamics for the chosen $J,\Delta,g_A,g_B,N$ and $\omega_E$.

\section{Ising coupling: canonical form of the time-dependent dephasing channel}
\label{sec:canonical}

The Schur structure of Eq.~\eqref{eq:schur} fixes the time-local
generator in closed form. Here we bring it to the canonical form of
Hall \emph{et al.}~\cite{Hall2014}, which uniquely fixes the canonical
rates. The associated operators retain phase and degeneracy freedoms,
as specified in Section~\ref{sec:ratesoperators}. Every result of this
section holds for an arbitrary fixed bulk preparation $\omega_E$ and
applies to every probe input $\rho_{AB}(0)$, pure or mixed, product or
entangled: the canonical form is a property of the channel $\Lambda_t$,
which enters only through the coherence factors $\Gamma_{ab}(t)$ of
Eq.~\eqref{eq:gamma}.

The notation distinguishes three quantities: $\Gamma_{ab}(t)$ is a
complex channel multiplier, $\kappa_{ab}(t)$ is its real logarithmic
decay rate, and $\gamma_k(t)$ is a real canonical rate. The pair of
configuration indices $a,b$ labels a coherence; the single index $k$
labels an eigenmode of the decoherence matrix. Both rates have units
of inverse time, whereas $\Gamma$ and $\Gamma_{ab}$ are dimensionless.

\subsection{Time-local generator of the Schur channel}
\label{sec:schurgenerator}

When the dynamical map is invertible, the reduced evolution obeys a
time-local master equation with generator
$\mathcal L_t=\dot\Lambda_t\Lambda_t^{-1}$~\cite{Hall2014}. For the
Schur channel this requires $\Gamma_{ab}(t)\ne0$ for all $a,b$
(Section~\ref{sec:invertibility}). Differentiating
$\rho_{ab}(t)=\Gamma_{ab}(t)r_{ab}$ and, for $\Gamma_{ab}(t)\ne0$,
substituting $r_{ab}=\rho_{ab}(t)/\Gamma_{ab}(t)$ gives
\begin{equation}
 \dot\rho_{ab}(t)=\eta_{ab}(t)\,\rho_{ab}(t),\qquad
 \eta_{ab}(t)=\frac{\dot\Gamma_{ab}(t)}{\Gamma_{ab}(t)}
 =\frac{d}{dt}\ln\Gamma_{ab}(t),
 \label{eq:schurgenerator}
\end{equation}
so $\mathcal L_t(X)_{ab}=\eta_{ab}(t)X_{ab}$, in agreement with
Eq.~\eqref{eq:generatorschur}. The constraints of
Eq.~\eqref{eq:correlation} imply $\eta_{aa}=0$ and
$\eta_{ba}=\eta_{ab}^*$, so the $4\times4$ matrix $\eta=(\eta_{ab})$ is
Hermitian with zero diagonal. All generator data are six complex
logarithmic derivatives. Writing each nonzero factor in polar form,
$\Gamma_{ab}=R_{ab}e^{i\theta_{ab}}$ with $R_{ab}>0$,
\begin{equation}
 \dot\Gamma_{ab}=(\dot R_{ab}+iR_{ab}\dot\theta_{ab})e^{i\theta_{ab}},\qquad
 \eta_{ab}=\frac{\dot R_{ab}}{R_{ab}}+i\dot\theta_{ab}.
\end{equation}
We therefore define, without assuming Markovianity,
\begin{equation}
 \eta_{ab}=-\kappa_{ab}+i\omega_{ab},\qquad
 \kappa_{ab}=-\frac{d}{dt}\ln\lvert\Gamma_{ab}\rvert,\qquad
 \omega_{ab}=\frac{d}{dt}\arg\Gamma_{ab}.
 \label{eq:kappaomega}
\end{equation}
Here $\kappa_{ab}=\kappa_{ba}$ and $\omega_{ab}=-\omega_{ba}$; the
derivative $\omega_{ab}$ is taken on a locally continuous phase branch.
Equivalently, $\omega_{ab}=\Im(\dot\Gamma_{ab}/\Gamma_{ab})$, which
avoids artificial jumps of the principal argument.
For any input with $r_{ab}\ne0$,
\begin{equation}
 \kappa_{ab}(t)=-\frac{1}{\lvert\rho_{ab}(t)\rvert}
 \frac{d\lvert\rho_{ab}(t)\rvert}{dt},
 \label{eq:fractionalrate}
\end{equation}
so $\kappa_{ab}$ is the fractional rate of loss of the coherence
between configurations $a$ and $b$, independent of the input amplitude.
It is an observed coherence damping rate, not in general an eigenvalue
of the canonical decoherence matrix.
Over an interval on which $\Gamma_{ab}$ does not vanish,
$\int_{t_1}^{t_2}\kappa_{ab}\,dt=\ln\bigl(\lvert\Gamma_{ab}(t_1)\rvert
/\lvert\Gamma_{ab}(t_2)\rvert\bigr)$.

\subsection{Canonical form}

Let $\{O_j\}_{j=1}^{15}$ be an orthonormal basis of traceless probe
operators, $\Tr O_j=0$ and $\Tr(O_j^\dagger O_k)=\delta_{jk}$. Any
time-local master equation of a $d$-dimensional system can be written
as~\cite{Hall2014}
\begin{equation}
 \dot\rho=-i[H(t),\rho]+\sum_{j,k}\dmat_{jk}(t)
 \Bigl(O_j\rho O_k^\dagger-\tfrac12\{O_k^\dagger O_j,\rho\}\Bigr),
 \label{eq:canonicalgeneral}
\end{equation}
with a Hermitian traceless $H(t)$ and a Hermitian decoherence matrix
$\dmat(t)$. Diagonalizing $\dmat$ gives
\begin{equation}
 \dot\rho=-i[H(t),\rho]+\sum_{k}\gamma_k(t)
 \Bigl(L_k(t)\rho L_k^\dagger(t)
 -\tfrac12\{L_k^\dagger(t)L_k(t),\rho\}\Bigr),
 \label{eq:canonicalform}
\end{equation}
where the $L_k(t)$ form an orthonormal basis of traceless operators.
The canonical decoherence rates $\gamma_k(t)$ are the eigenvalues of
$\dmat(t)$; they are uniquely determined, invariant under the choice
of the basis $\{O_j\}$, time dependent, and can be
negative~\cite{Hall2014}. The Hamiltonian $H(t)$ acts on the probe pair
and is distinct from the total probe--bulk Hamiltonian $H^{(\ZZ)}$. We
construct this form explicitly for $\mathcal L_t$ of
Eq.~\eqref{eq:schurgenerator}.

\subsection{Diagonal operator basis}
\label{sec:diagonalsubspace}

Define the three diagonal, traceless and Hermitian operators
\begin{equation}
 F_1=\frac{\sigma_A^z\otimes I}{2},\qquad
 F_2=\frac{I\otimes\sigma_B^z}{2},\qquad
 F_3=\frac{\sigma_A^z\otimes\sigma_B^z}{2}.
 \label{eq:diagonalbasis}
\end{equation}
Each is a product of Paulis divided by two, so
$\Tr F_\mu=0$ and $\Tr(F_\mu F_\nu)=\delta_{\mu\nu}$. For the configuration
$a=(s_A^a,s_B^a)$, $F_\mu\ket a=g_\mu^a\ket a$ with
$\mathbf g^a=(g_1^a,g_2^a,g_3^a)=(s_A^a,\,s_B^a,\,2s_A^as_B^a)$. Collecting
these eigenvalues into the $4\times3$ matrix $G=(g_\mu^a)_{a\mu}$,
\begin{equation}
 G=\frac12\begin{pmatrix}
 1&1&1\\ 1&-1&-1\\ -1&1&-1\\ -1&-1&1
 \end{pmatrix},
 \qquad
 \boxed{G^{\mathsf T}G=I_3,\qquad
 G^{\mathsf T}\mathbf1=0,\qquad
 GG^{\mathsf T}=I_4-\tfrac14\mathbf1\mathbf1^{\mathsf T},}
 \label{eq:Gmatrix}
\end{equation}
where $\mathbf1=(1,1,1,1)^{\mathsf T}$. The first two identities follow
by direct multiplication: each column has four entries $\pm\frac12$,
distinct columns differ in sign at exactly two rows, and each column has
two entries of each sign. The third follows because the three columns
and $\mathbf1/2$ form an orthonormal basis of $\mathbb R^4$. Row $a$ of
$G$, viewed as a column vector, is $\mathbf g^a$; hence
$\sum_a\mathbf g^a=G^{\mathsf T}\mathbf1=0$.

Complete $\{F_1,F_2,F_3\}$ with any twelve orthonormal traceless
off-diagonal operators to obtain a basis $\{O_j\}$. We take the ansatz
of a diagonal Hamiltonian
\begin{equation}
 \Heff=\sum_ah_aP_a,\qquad h_a\in\mathbb R,\qquad \sum_ah_a=0,
\end{equation}
where the last condition fixes the irrelevant scalar gauge, and a
decoherence matrix supported on the diagonal operators,
\begin{equation}
 \dmat=S+iA,\qquad S=S^{\mathsf T}\in\mathbb R^{3\times3},\qquad
 A=-A^{\mathsf T}\in\mathbb R^{3\times3},
 \label{eq:dsplit}
\end{equation}
indexed by $\mu,\nu\in\{1,2,3\}$, with all other entries zero. Any
Hermitian $3\times3$ matrix has this decomposition, with nine real
parameters. \emph{No positivity assumption} is imposed on $\dmat$ or its
eigenvalues: these may have either sign while $\Lambda_t$ remains CPTP.
The ansatz reads
\begin{equation}
 \mathcal L_t[X]=-i[\Heff,X]+
 \sum_{\mu,\nu=1}^{3}\dmat_{\mu\nu}
 \Bigl(F_\mu XF_\nu-\tfrac12\{F_\nu F_\mu,X\}\Bigr).
 \label{eq:ansatz}
\end{equation}

\subsection{Action of each term on a matrix unit}

Because every $F_\mu$ and $\Heff$ is diagonal,
\begin{equation}
 F_\mu E_{ab}F_\nu=g_\mu^ag_\nu^bE_{ab},\qquad
 F_\nu F_\mu E_{ab}=g_\nu^ag_\mu^aE_{ab},\qquad
 E_{ab}F_\nu F_\mu=g_\nu^bg_\mu^bE_{ab},
\end{equation}
\begin{equation}
 [\Heff,E_{ab}]=(h_a-h_b)E_{ab}.
\end{equation}
Thus the ansatz maps $E_{ab}$ to a multiple of itself, with the exact
scalar coefficient
\begin{equation}
 \eta_{ab}=-i(h_a-h_b)
 +\sum_{\mu\nu}\dmat_{\mu\nu}g_\mu^ag_\nu^b
 -\frac12\sum_{\mu\nu}\dmat_{\mu\nu}(g_\nu^ag_\mu^a+g_\nu^bg_\mu^b).
\end{equation}
In vector notation,
\begin{equation}
 \boxed{\eta_{ab}=-i(h_a-h_b)
 +\mathbf g^{a\mathsf T}\dmat\,\mathbf g^b
 -\tfrac12\mathbf g^{a\mathsf T}\dmat\,\mathbf g^a
 -\tfrac12\mathbf g^{b\mathsf T}\dmat\,\mathbf g^b.}
 \label{eq:elementwise}
\end{equation}
This must equal $\dot\Gamma_{ab}/\Gamma_{ab}$ for every $a,b$.
Because the $\mathbf g^a$ are real and $A$ is real antisymmetric,
$\mathbf x^{\mathsf T}A\mathbf x=0$ for every real $\mathbf x$, so
$\mathbf g^{a\mathsf T}\dmat\mathbf g^a=\mathbf g^{a\mathsf T}S\mathbf g^a$ is real
and $\mathbf g^{a\mathsf T}\dmat\mathbf g^b
=\mathbf g^{a\mathsf T}S\mathbf g^b+i\,\mathbf g^{a\mathsf T}A\mathbf g^b$.
Separating real and imaginary parts, with
$\Delta_{ab}=\mathbf g^a-\mathbf g^b$,
\begin{align}
 \Re\eta_{ab}&=\mathbf g^{a\mathsf T}S\mathbf g^b
 -\tfrac12\mathbf g^{a\mathsf T}S\mathbf g^a-\tfrac12\mathbf g^{b\mathsf T}S\mathbf g^b
 =-\tfrac12\Delta_{ab}^{\mathsf T}S\,\Delta_{ab},\\
 \Im\eta_{ab}&=-(h_a-h_b)+\mathbf g^{a\mathsf T}A\,\mathbf g^b .
\end{align}
Comparing with Eq.~\eqref{eq:kappaomega},
\begin{align}
 \kappa_{ab}&=\tfrac12\Delta_{ab}^{\mathsf T}S\,\Delta_{ab},
 \label{eq:kappaS}\\
 \omega_{ab}&=-(h_a-h_b)+\mathbf g^{a\mathsf T}A\,\mathbf g^b.
 \label{eq:omegaA}
\end{align}
The symmetric part governs instantaneous damping; the antisymmetric
part contributes phases beyond those of a diagonal Hamiltonian.

\subsection{Matrix form and unique inversion}
\label{sec:inversion}

Define $q_a=\frac12\mathbf g^{a\mathsf T}S\mathbf g^a$ and the vectors
$h=(h_a)_a$, $q=(q_a)_a$. Equation~\eqref{eq:elementwise} for all sixteen
entries, including $a=b$ where it gives $\eta_{aa}=0$, is the matrix
identity
\begin{equation}
 \eta=G\dmat G^{\mathsf T}+(-ih-q)\mathbf1^{\mathsf T}+\mathbf1(ih-q)^{\mathsf T},
 \label{eq:etamatrix}
\end{equation}
with ordinary transposes: the $(a,b)$ entry of the right side is
$\mathbf g^{a\mathsf T}\dmat\mathbf g^b-ih_a-q_a+ih_b-q_b$. Multiplying on
the left by $G^{\mathsf T}$ and on the right by $G$, and using
Eq.~\eqref{eq:Gmatrix},
\begin{equation}
 G^{\mathsf T}\eta G=(G^{\mathsf T}G)\dmat(G^{\mathsf T}G)
 +G^{\mathsf T}(-ih-q)(\mathbf1^{\mathsf T}G)
 +(G^{\mathsf T}\mathbf1)(ih-q)^{\mathsf T}G=\dmat,
\end{equation}
that is,
\begin{equation}
 \boxed{\dmat(t)=G^{\mathsf T}\eta(t)G,\qquad
 \dmat_{\mu\nu}=\sum_{a,b}g_\mu^a\,\eta_{ab}\,g_\nu^b,\qquad
 S=\Re\dmat,\quad A=\Im\dmat.}
 \label{eq:dfrometa}
\end{equation}
This reconstructs all nine real parameters of $\dmat$. Its antisymmetric
part is $A_{\mu\nu}=\sum_{a,b}g_\mu^a\omega_{ab}g_\nu^b$, since $g$ is real.

For the Hamiltonian, sum Eq.~\eqref{eq:omegaA} over $b$ at fixed $a$:
\begin{equation}
 \sum_b\omega_{ab}=-4h_a+\sum_bh_b
 +\mathbf g^{a\mathsf T}A\sum_b\mathbf g^b=-4h_a,
\end{equation}
using $\sum_bh_b=0$ and $\sum_b\mathbf g^b=0$. Therefore
\begin{equation}
 \boxed{h_a(t)=-\frac14\sum_{b}\omega_{ab}(t)
 =-\frac14\sum_{b}\Im\eta_{ab}(t),\qquad
 \Heff(t)=\sum_ah_a(t)P_a .}
 \label{eq:canonicalH}
\end{equation}
The four diagonal traceless matrices are spanned by $F_1,F_2,F_3$, so
$\Heff$ may contain $\sigma_A^z$, $\sigma_B^z$ and an induced
$\sigma_A^z\sigma_B^z$ term.

\paragraph{Uniqueness.}
The map $(\dmat,h)\mapsto\eta$ of Eq.~\eqref{eq:etamatrix}, with $\dmat$
Hermitian $3\times3$ (nine real parameters) and $h$ real with
$\sum_ah_a=0$ (three real parameters), is real-linear and takes values
in the Hermitian $4\times4$ matrices with zero diagonal (six complex, or
twelve real, parameters). It is injective: if $\eta=0$, then
Eq.~\eqref{eq:dfrometa} gives $\dmat=0$, and Eq.~\eqref{eq:canonicalH}
gives $h=0$. An injective linear map between real spaces of equal
dimension twelve is bijective. Hence for every Hermitian $\eta$ with zero
diagonal there is exactly one pair $(\dmat,\Heff)$ of the form
\eqref{eq:ansatz}, given by Eqs.~\eqref{eq:dfrometa}
and~\eqref{eq:canonicalH}. The ansatz is of the general form
\eqref{eq:canonicalgeneral}, and uniqueness of the canonical
form~\cite{Hall2014} identifies it as the canonical form of $\mathcal L_t$.
Therefore:
\begin{proposition}\label{prop:threerates}
At every time at which $\Lambda_t$ is invertible, the canonical
Hamiltonian of the Ising channel is diagonal in the configuration basis,
and its decoherence matrix is supported on
$\operatorname{span}\{F_1,F_2,F_3\}$. At most three canonical rates are
nonzero; the remaining twelve vanish identically.
\end{proposition}
No off-diagonal operator, and hence no operator that transfers
population between configurations, appears in the generator. This is
consistent with the fixed populations, $\Gamma_{aa}(t)=1$.

\subsection{Dissipative part: the coherence magnitudes}
\label{sec:dissipative}

The six difference vectors and Eq.~\eqref{eq:kappaS} give
\begin{center}
\begin{tabular}{@{}ccl@{}}
\toprule
pair $ab$&$\Delta_{ab}$&$\kappa_{ab}=\tfrac12\Delta_{ab}^{\mathsf T}S\Delta_{ab}$\\\midrule
$00,01$&$(0,1,1)$&$\tfrac12(S_{22}+S_{33}+2S_{23})$\\
$10,11$&$(0,1,-1)$&$\tfrac12(S_{22}+S_{33}-2S_{23})$\\
$00,10$&$(1,0,1)$&$\tfrac12(S_{11}+S_{33}+2S_{13})$\\
$01,11$&$(1,0,-1)$&$\tfrac12(S_{11}+S_{33}-2S_{13})$\\
$00,11$&$(1,1,0)$&$\tfrac12(S_{11}+S_{22}+2S_{12})$\\
$01,10$&$(1,-1,0)$&$\tfrac12(S_{11}+S_{22}-2S_{12})$\\\bottomrule
\end{tabular}
\end{center}
This linear map from the six entries of $S$ to the six $\kappa_{ab}$ is
invertible. Differences within each row pair give the off-diagonal
entries; with the pair sums
\begin{align}
 K_A&=\kappa_{00,10}+\kappa_{01,11}=S_{11}+S_{33},\qquad
 K_B=\kappa_{00,01}+\kappa_{10,11}=S_{22}+S_{33},\notag\\
 K_{AB}&=\kappa_{00,11}+\kappa_{01,10}=S_{11}+S_{22},
\end{align}
for coherences in which only $A$, only $B$, or both probes flip, the
inverse is
\begin{align}
 S_{11}&=\tfrac12(K_A+K_{AB}-K_B),&
 S_{12}&=\tfrac12(\kappa_{00,11}-\kappa_{01,10}),\notag\\
 S_{22}&=\tfrac12(K_B+K_{AB}-K_A),&
 S_{13}&=\tfrac12(\kappa_{00,10}-\kappa_{01,11}),
 \label{eq:Sinverse}\\
 S_{33}&=\tfrac12(K_A+K_B-K_{AB}),&
 S_{23}&=\tfrac12(\kappa_{00,01}-\kappa_{10,11}).\notag
\end{align}
This is an explicit alternative to $S=\Re(G^{\mathsf T}\eta G)$.
The entries have a direct reading:
\begin{itemize}
\item $S_{11}$ and $S_{22}$ describe dephasing of each probe along $z$.
\item $S_{12}>0$ when the coherence between $00$ and $11$ decays faster
than that between $01$ and $10$, as for noise acting equally on both
probes; $S_{12}<0$ in the opposite, anticorrelated case.
\item $S_{13}$ ($S_{23}$) measures how the dephasing of $A$ ($B$)
depends on the configuration of the other probe.
\item $S_{33}$ is the diagonal decoherence-matrix coefficient associated
with $F_3=\sigma_A^z\sigma_B^z/2$.
\end{itemize}
These entries are coefficients in the chosen operator basis, rather
than canonical rates. They coincide with canonical rates when the full
matrix $\dmat=S+iA$ is diagonal in that basis
(Section~\ref{sec:ratesoperators}).
If $\Gamma$ factorizes into two independent single-probe dephasing
maps, $\Gamma_{(s_As_B),(s_A's_B')}=\alpha_{s_As_A'}\beta_{s_Bs_B'}$,
then $\kappa_{00,10}=\kappa_{01,11}=\kappa^A$,
$\kappa_{00,01}=\kappa_{10,11}=\kappa^B$ and
$\kappa_{00,11}=\kappa_{01,10}=\kappa^A+\kappa^B$, so
Eq.~\eqref{eq:Sinverse} gives $S=\diag(2\kappa^A,2\kappa^B,0)$. Nonzero
$S_{12}$, $S_{13}$, $S_{23}$ or $S_{33}$ therefore signal dephasing
correlated through the common bulk.

For $a\ne b$, each $\Delta_{ab}$ has squared norm $2$. With the unit vector
$\mathbf e_{ab}=\Delta_{ab}/\sqrt2$ and the normalized operator
$F_{ab}=\sum_\mu(\mathbf e_{ab})_\mu F_\mu$,
\begin{equation}
 \kappa_{ab}=\mathbf e_{ab}^{\mathsf T}S\,\mathbf e_{ab}
 =\mathbf e_{ab}^{\mathsf T}\dmat\,\mathbf e_{ab},
 \label{eq:kappaprojection}
\end{equation}
since $\mathbf e^{\mathsf T}A\mathbf e=0$. Each coherence decay rate is
thus the decoherence matrix projected onto one fixed traceless
diagonal operator; for example,
$F_{01,10}=(\sigma_A^z-\sigma_B^z)/(2\sqrt2)$ and
$F_{00,11}=(\sigma_A^z+\sigma_B^z)/(2\sqrt2)$.

\subsection{Coherent part: the coherence phases}
\label{sec:coherent}

The energies $h_a$ cancel from a sum of phase rates around a closed loop
of three configurations. For an oriented triple $(a,b,c)$,
Eq.~\eqref{eq:omegaA} gives
\begin{equation}
 \mathcal C_{abc}\equiv\omega_{ab}+\omega_{bc}+\omega_{ca}
 =\mathbf g^{a\mathsf T}A\mathbf g^b+\mathbf g^{b\mathsf T}A\mathbf g^c
 +\mathbf g^{c\mathsf T}A\mathbf g^a,
\end{equation}
because $(h_a-h_b)+(h_b-h_c)+(h_c-h_a)=0$.
Write $A$ through a vector $\mathbf a=(a_1,a_2,a_3)$,
$A_{\mu\nu}=\epsilon_{\mu\nu\lambda}a_\lambda$, so that
$\mathbf x^{\mathsf T}A\mathbf y=\mathbf a\cdot(\mathbf x\times\mathbf y)$.
Since $(\mathbf g^b-\mathbf g^a)\times(\mathbf g^c-\mathbf g^a)
=\mathbf g^a\times\mathbf g^b+\mathbf g^b\times\mathbf g^c+\mathbf g^c\times\mathbf g^a$,
\begin{equation}
 \mathcal C_{abc}
 =\mathbf a\cdot\bigl[(\mathbf g^b-\mathbf g^a)\times
 (\mathbf g^c-\mathbf g^a)\bigr].
 \label{eq:loopsum}
\end{equation}
The four configurations have three independent loops. With
$\mathbf g^{01}-\mathbf g^{00}=(0,-1,-1)$, $\mathbf g^{10}-\mathbf g^{00}=(-1,0,-1)$
and $\mathbf g^{11}-\mathbf g^{00}=(-1,-1,0)$, the cross products are
$(1,1,-1)$, $(-1,1,-1)$ and $(-1,1,1)$, respectively, for the loops
\begin{align}
 \mathcal C_1&\equiv \mathcal C_{00,01,10}=a_1+a_2-a_3,\notag\\
 \mathcal C_2&\equiv \mathcal C_{00,01,11}=-a_1+a_2-a_3,\\
 \mathcal C_3&\equiv \mathcal C_{00,10,11}=-a_1+a_2+a_3,\notag
\end{align}
which determine
\begin{equation}
 a_1=\tfrac12(\mathcal C_1-\mathcal C_2),\qquad
 a_2=\tfrac12(\mathcal C_1+\mathcal C_3),\qquad
 a_3=\tfrac12(\mathcal C_3-\mathcal C_2).
 \label{eq:Asolution}
\end{equation}
This agrees with $A=\Im(G^{\mathsf T}\eta G)$ and makes its physical
content explicit: the matrix $A$ carries the part of the phase evolution
that cannot be absorbed into four configuration energies, and it
vanishes iff all loop sums $\mathcal C_{abc}$ vanish. With $A$ fixed,
the remaining phase rates $\zeta_{ab}=-\omega_{ab}+\mathbf g^{a\mathsf T}A\,\mathbf g^b$
equal $h_a-h_b$ and reproduce Eq.~\eqref{eq:canonicalH}, since
$\frac14\sum_b\zeta_{ab}=-\frac14\sum_b\omega_{ab}$. The Hamiltonian is a
configuration-dependent energy shift of the probes induced by the bulk.
For a factorized $\Gamma$ (previous subsection) every loop sum vanishes,
because each phase rate is a sum of single-probe phase rates that cancel
around a closed loop; hence $A=0$ in that case.

\subsection{Canonical rates and Lindblad operators}
\label{sec:ratesoperators}

Diagonalize the Hermitian matrix of Eq.~\eqref{eq:dfrometa} at each
regular time,
\begin{equation}
 \dmat(t)=V(t)\diag(\gamma_1(t),\gamma_2(t),\gamma_3(t))V(t)^\dagger,
 \qquad V^\dagger V=I_3,
\end{equation}
so that
\begin{equation}
 \boxed{\{\gamma_k(t)\}_{k=1}^3=\spec\bigl[S(t)+iA(t)\bigr],\qquad
 \gamma_{4},\ldots,\gamma_{15}=0.}
 \label{eq:canonicalrates}
\end{equation}
We order these three rates as $\gamma_1\le\gamma_2\le\gamma_3$. The
additional twelve zeros are listed separately, so the indices
$1,\ldots,15$ do not constitute a globally sorted list. Define
\begin{equation}
 \boxed{L_k(t)=\sum_{\mu=1}^3V_{\mu k}(t)F_\mu
 =\frac12\bigl[V_{1k}\sigma_A^z+V_{2k}\sigma_B^z
 +V_{3k}\sigma_A^z\sigma_B^z\bigr].}
 \label{eq:lindbladops}
\end{equation}
Substituting $\dmat_{\mu\nu}=\sum_kV_{\mu k}\gamma_kV_{\nu k}^*$ into
Eq.~\eqref{eq:ansatz} and using $F_\nu^\dagger=F_\nu$,
\begin{equation}
 \sum_{\mu\nu}\dmat_{\mu\nu}F_\mu XF_\nu
 =\sum_k\gamma_k\Bigl(\sum_\mu V_{\mu k}F_\mu\Bigr)X
 \Bigl(\sum_\nu V_{\nu k}F_\nu\Bigr)^\dagger=\sum_k\gamma_kL_kXL_k^\dagger,
\end{equation}
and likewise $\sum_{\mu\nu}\dmat_{\mu\nu}F_\nu F_\mu=\sum_k\gamma_kL_k^\dagger L_k$.
The $L_k$ are traceless and satisfy
$\Tr(L_j^\dagger L_k)=\sum_{\mu\nu}V_{\mu j}^*V_{\nu k}\delta_{\mu\nu}=\delta_{jk}$.
They are diagonal in the configuration basis,
\begin{equation}
 L_k(t)\ket a=\Bigl(\sum_\mu V_{\mu k}(t)g^a_\mu\Bigr)\ket a,
 \label{eq:lindbladaction}
\end{equation}
so each $L_k$ imprints a configuration-dependent amplitude and phase.
If $A=0$, the eigenvectors can be chosen real and the $L_k$ are
Hermitian; for $A\ne0$ they are diagonal but generally non-Hermitian.
For nondegenerate rates each $L_k$ is unique up to a phase; within a
degenerate rate subspace, orthonormal eigenoperators may be mixed by
a unitary transformation~\cite{Hall2014}. Ordering rates at each time
does not track an eigenoperator continuously through a crossing.

The canonical master equation of the Ising channel is therefore
\begin{equation}
 \dot\rho=-i[\Heff(t),\rho]+\sum_{k=1}^3\gamma_k(t)
 \Bigl(L_k(t)\rho L_k^\dagger(t)
 -\tfrac12\{L_k^\dagger(t)L_k(t),\rho\}\Bigr),
 \label{eq:canonicalmaster}
\end{equation}
with $\Heff$, $\gamma_k$ and $L_k$ determined by $\Gamma(t)$ alone. All
operators in it are diagonal and mutually commute, so it fixes the
populations and is unital. The $L_k$ are not raising or lowering
operators: such operators would move population, which
Proposition~\ref{prop:threerates} excludes for this channel.

\subsection{The generator at the initial time}
\label{sec:initial}

By Eq.~\eqref{eq:initialderivative},
$\eta_{ab}(0)=\dot\Gamma_{ab}(0)=-i(\varepsilon_a-\varepsilon_b)$.
Hence $\kappa_{ab}(0)=0$ for every pair, and the invertible map of
Eq.~\eqref{eq:Sinverse} gives $S(0)=0$. In matrix form,
$\eta(0)=-i\varepsilon\mathbf1^{\mathsf T}+i\mathbf1\varepsilon^{\mathsf T}$
with $\varepsilon=(\varepsilon_a)_a$, so Eq.~\eqref{eq:dfrometa} and
$\mathbf1^{\mathsf T}G=0$ give
\begin{equation}
 \dmat(0)=G^{\mathsf T}\eta(0)G=0,\qquad
 h_a(0)=\varepsilon_a-\frac14\sum_b\varepsilon_b .
\end{equation}
All canonical rates vanish at $t=0$, and the initial motion is purely
coherent, generated by the bulk-averaged conditional energies
$\Tr_E(\omega_EH_a)$. This is the general behaviour of a product initial
state under a finite-dimensional unitary dilation; dissipative
contributions to $\dmat$ appear only at $t>0$.

\subsection{Sign of the rates and CP divisibility}
\label{sec:ratesign}

Following Hall \emph{et al.}~\cite{Hall2014}, a time-local master
equation is called Markovian at time $t$ iff all canonical rates are
nonnegative, and non-Markovian iff at least one rate is strictly
negative. On an interval where the map is differentiable and invertible,
nonnegative rates throughout the interval are equivalent to CP
divisibility~\cite{Hall2014}:
\begin{equation}
 \text{CP-divisible}\iff\dmat(t)\succeq0\iff\gamma_k(t)\ge0\quad\forall k
 \ \text{(almost everywhere)}.
\end{equation}
A negative rate does not invalidate Eq.~\eqref{eq:canonicalmaster}, and
$\Lambda_t$ remains CPTP for every $t\ge0$, because $\Gamma(t)\succeq0$
with unit diagonal; it signals that the intermediate maps near $t$ are
not completely positive, consistent with the finite-interval test of
Eq.~\eqref{eq:divisibility}.

For the Ising channel the sign structure follows from $S$ and $A$:
write $\gamma_{\min}=\gamma_1$ and $\gamma_{\max}=\gamma_3$ for the
extreme eigenvalues of the $3\times3$ diagonal-operator block.
\begin{enumerate}
\item The smallest rate obeys $\gamma_{\min}\le\lambda_{\min}(S)$.
Indeed $S=(\dmat+\dmat^{\mathsf T})/2$, and $\dmat^{\mathsf T}=\dmat^*$
has the same spectrum as $\dmat$, so Weyl's inequality gives
$\lambda_{\min}(S)\ge\frac12[\lambda_{\min}(\dmat)+\lambda_{\min}(\dmat^{\mathsf T})]
=\gamma_{\min}$. A non-positive-semidefinite $S$
therefore implies a negative canonical rate; the phase part $A$ cannot
restore positivity.
\item By Eq.~\eqref{eq:kappaprojection} and the Rayleigh-quotient bounds
for a Hermitian matrix and a unit vector, for every distinct pair $a\ne b$,
\begin{equation}
 \gamma_{\min}(t)\le\kappa_{ab}(t)\le\gamma_{\max}(t).
 \label{eq:ratebounds}
\end{equation}
Hence a negative rate occurs whenever any coherence magnitude
$\lvert\Gamma_{ab}\rvert$ increases. The converse fails: a negative rate
can occur with every $\kappa_{ab}$ positive. Likewise
$\gamma_{\min}\le S_{\mu\mu}=\mathbf u_\mu^{\mathsf T}\dmat\mathbf u_\mu$ for each
unit vector $\mathbf u_\mu$; in particular $K_{AB}>K_A+K_B$, i.e.\
$S_{33}<0$, implies a negative rate even if every $\kappa_{ab}$ is
positive.
\item $S\succeq0$ iff there is a real $3\times3$ matrix $M$ with
$S=M^{\mathsf T}M$, i.e.\
$\kappa_{ab}=\frac12\lvert M\Delta_{ab}\rvert^2$ for all pairs.
In that case, the six decay rates are one half of squared Euclidean
distances between four points $M\mathbf g^a$ representing the
configurations.
\end{enumerate}

The complete condition is $\dmat=S+iA\succeq0$. Although
$S\succeq0$ and nonnegative $\kappa_{ab}$ are necessary for this
condition, neither alone is sufficient. For example,
$S=I_3$ with $A_{12}=-A_{21}=2$ and all other entries of $A$ zero
has $\kappa_{ab}=1$ for all $a\ne b$, but
$\spec(S+iA)=\{1-2,1,1+2\}=\{-1,1,3\}$.

\subsection{Singular times}

The rates are built from the logarithmic derivatives
$\eta_{ab}=\dot\Gamma_{ab}/\Gamma_{ab}$. The inverse-map construction
fails when a coherence factor vanishes~\cite{Hall2014}. For the finite
system with a time-independent Hamiltonian, each $\Gamma_{ab}(t)$ is
analytic in $t$ (a finite sum of exponentials). Near a zero of order
$r\ge1$ at $t_*$,
\begin{equation}
 \Gamma_{ab}(t)=c(t-t_*)^r[1+O(t-t_*)],\quad c\ne0,
 \qquad \eta_{ab}(t)=\frac{r}{t-t_*}+O(1).
 \label{eq:singularfactor}
\end{equation}
The second relation follows from
$\eta_{ab}=\frac{d}{dt}\ln\Gamma_{ab}=r/(t-t_*)+\frac{d}{dt}\ln[1+O(t-t_*)]$,
where the last term is analytic. Since the pole is real,
$\kappa_{ab}=-r/(t-t_*)+O(1)$ diverges, whereas
$\omega_{ab}=\Im\eta_{ab}=O(1)$ remains bounded. For odd $r$ the phase
$\arg\Gamma_{ab}$ jumps by $\pi$ across $t_*$, because $(t-t_*)^r$ changes
sign. By Eq.~\eqref{eq:ratebounds}, at least one canonical rate becomes
unbounded: $\gamma_{\max}\to+\infty$ as $t\to t_*^-$ and
$\gamma_{\min}\to-\infty$ as $t\to t_*^+$. At $t_*$ the rates defined by
$\dot\Lambda_t\Lambda_t^{-1}$ are undefined. The time-local form
applies on neighboring intervals where all coherence factors are nonzero.
The state and its derivative remain bounded:
$\rho_{ab}(t)=\Gamma_{ab}(t)r_{ab}$ and
$\lvert\dot\rho_{ab}\rvert\le\lvert\dot\Gamma_{ab}\rvert\lvert r_{ab}\rvert$.
According to Eq.~\eqref{eq:fractionalrate}, $\kappa_{ab}$ measures a
fractional rather than an absolute change. Its signed integral is finite
on a closed interval containing no zeros, as shown below
Eq.~\eqref{eq:fractionalrate}. At an exact zero the corresponding
one-sided logarithmic integrals diverge; a principal-value cancellation
does not define a finite ordinary integral across the singularity.
Lemma~\ref{lem:nosingular} shows that this divergence is unavoidable: an
invertible map cannot reach a singular one under a locally integrable
generator.

\subsection{Computational procedure and consistency tests}
\label{sec:procedure}

For each chosen $(N,J,\Delta,g_A,g_B,\omega_E)$ and time $t$:
\begin{enumerate}
\item Build the four bulk matrices $H_a$ of Eq.~\eqref{eq:conditional}
and propagate $U_a(t)$, or propagate the vectors of a decomposition of
$\omega_E$.
\item Compute the six complex values
$\Gamma_{ab}=\Tr(\omega_EU_b^\dagger U_a)$, $a<b$
[Eq.~\eqref{eq:gamma}].
\item Compute $\dot\Gamma_{ab}=-i\Tr[\omega_EU_b^\dagger(H_a-H_b)U_a]$
[Eq.~\eqref{eq:gammadot}].
\item If each $\Gamma_{ab}\ne0$, set $\eta_{ab}=\dot\Gamma_{ab}/\Gamma_{ab}$,
$\eta_{ba}=\eta_{ab}^*$ and $\eta_{aa}=0$.
\item Form $\dmat=G^{\mathsf T}\eta G$ and $h_a=-\frac14\sum_b\Im\eta_{ab}$
[Eqs.~\eqref{eq:dfrometa} and~\eqref{eq:canonicalH}].
\item Diagonalize only the $3\times3$ Hermitian matrix $\dmat$, obtaining
$\gamma_k$ and the diagonal operators $L_k$ [Eq.~\eqref{eq:lindbladops}].
\end{enumerate}
No $16\times16$ diagonalization is required. At every regular time the
calculation can be checked independently:
\begin{enumerate}
\item $\Gamma=\Gamma^\dagger\succeq0$, $\Gamma_{aa}=1$;
$\eta=\eta^\dagger$, $\eta_{aa}=0$.
\item $\dmat=\dmat^\dagger$; $\sum_ah_a=0$;
$\Tr(L_j^\dagger L_k)=\delta_{jk}$.
\item Reconstruct all six complex multipliers from the canonical form,
$\eta^{\mathrm{rec}}_{ab}$ given by the right side of
Eq.~\eqref{eq:elementwise}, and check
$\eta^{\mathrm{rec}}_{ab}=\dot\Gamma_{ab}/\Gamma_{ab}$.
\item Apply $\mathcal L_t$ to all sixteen $E_{ab}$ and check
$\mathcal L_t(E_{ab})=\eta_{ab}E_{ab}$.
\end{enumerate}
These tests verify the \emph{entire two-probe channel generator}, not
merely its action on the special input $\ket{++}$. Step~3 can also be
replaced by a finite difference of $\Gamma$; the analytic form avoids
the associated discretization error.

\part{Exchange coupling}
\label{part:exchange}

\section{Exchange coupling: exact symmetry reduction}
\label{sec:ex}

\subsection{Total charge and Hamiltonian sectors}

For the field-free model define
\begin{equation}
 M_S=S_A^z+S_B^z,\qquad M_E=\sum_{j=1}^N S_j^z,
 \qquad M=M_S+M_E.
\end{equation}
The bulk XXZ terms conserve $M_E$, and each exchange term
$S_A^+S_1^-+S_A^-S_1^+$ (or its right-contact analogue) transfers one unit
between a probe and the bulk while conserving their sum. Hence
$[H^{(\ex)},M]=0$, while individual probe $S^z$ values generally are not
conserved. The total-charge sector $m$ contains the basis states with
$(N+2)/2-m$ down spins, so
\begin{equation}
 \dim\HH_m=\binom{N+2}{(N+2)/2-m}.
\end{equation}
This Hamiltonian decomposition requires no assumption on the input;
coherences between occupied sectors must still be retained.
There is no four-branch controlled propagator of the form
\eqref{eq:controlled} for exchange coupling.

\subsection{Bulk assumption and proof of phase covariance}

For the channel reduction assume
\begin{equation}
 [\omega_E,M_E]=0.
 \label{eq:bulksym}
\end{equation}
Support in a single well-defined bulk magnetization sector is sufficient.
An incoherent mixture of sectors also suffices. A coherent superposition
of different bulk magnetizations generally does not.

Let $V_\theta=e^{-i\theta M_S}$ and $W_\theta=e^{-i\theta M_E}$, so that
$V_\theta\otimes W_\theta=e^{-i\theta M}$. Conservation gives
$[U_t,V_\theta\otimes W_\theta]=0$, while Eq.~\eqref{eq:bulksym} gives
$W_\theta\omega_EW_\theta^\dagger=\omega_E$. Then
\begin{align}
 \Lambda_t(V_\theta XV_\theta^\dagger)
 &=\Tr_E\bigl[U_t(V_\theta\otimes W_\theta)(X\otimes\omega_E)
 (V_\theta\otimes W_\theta)^\dagger U_t^\dagger\bigr]\notag\\
 &=\Tr_E\bigl[(V_\theta\otimes W_\theta)U_t(X\otimes\omega_E)
 U_t^\dagger(V_\theta\otimes W_\theta)^\dagger\bigr]\notag\\
 &=V_\theta\,\Tr_E\bigl[(I\otimes W_\theta)U_t(X\otimes\omega_E)U_t^\dagger
 (I\otimes W_\theta^\dagger)\bigr]\,V_\theta^\dagger
 =V_\theta\Lambda_t(X)V_\theta^\dagger,
\end{align}
where the last step uses cyclicity of the partial trace for operators
acting only on $E$. Hence
\begin{equation}
 \boxed{\Lambda_t(V_\theta XV_\theta^\dagger)
       =V_\theta\Lambda_t(X)V_\theta^\dagger.}
 \label{eq:covariance}
\end{equation}
This is $U(1)$ phase covariance. It permits exchange of magnetization
with the bulk; it does not imply $\Lambda_t^*(M_S)=M_S$ or unitality.

\subsection{Five superoperator blocks}

The probe charges are $(m_{00},m_{01},m_{10},m_{11})=(1,0,0,-1)$.
The operator $E_{ab}$ has coherence order $q=m_a-m_b$, meaning
$[M_S,E_{ab}]=qE_{ab}$ and $V_\theta E_{ab}V_\theta^\dagger=e^{-iq\theta}E_{ab}$.
Write $\Lambda_t(E_{ab})=\sum_{c,d}c_{cd}E_{cd}$. Covariance gives
\begin{equation}
 \sum_{c,d}e^{-i(m_c-m_d)\theta}c_{cd}E_{cd}
 =V_\theta\Lambda_t(E_{ab})V_\theta^\dagger
 =\Lambda_t(V_\theta E_{ab}V_\theta^\dagger)
 =e^{-iq\theta}\sum_{c,d}c_{cd}E_{cd}
\end{equation}
for every $\theta$. Hence $c_{cd}=0$ unless $m_c-m_d=q$, and
\begin{equation}
 \Lambda_t(\mathcal V_q)\subseteq\mathcal V_q,\qquad
 \mathcal V_q=\operatorname{span}\{E_{ab}:m_a-m_b=q\}.
\end{equation}
The ordered bases $\mathcal B_q$ of these spaces are
\begin{center}
\begin{tabular}{@{}ccl@{}}
\toprule
$q$&Dimension&Basis of $\mathcal V_q$\\\midrule
$0$&6&$E_{00,00},E_{01,01},E_{10,10},E_{11,11},E_{01,10},E_{10,01}$\\
$+1$&4&$E_{00,01},E_{00,10},E_{01,11},E_{10,11}$\\
$-1$&4&Adjoints of the $q=+1$ basis\\
$+2$&1&$E_{00,11}$\\
$-2$&1&$E_{11,00}$\\\bottomrule
\end{tabular}
\end{center}
Thus the $16\times16$ superoperator has the exact reduction
\begin{equation}
 \boxed{\Lambda_t\simeq T^{(-2)}\oplus T^{(-1)}\oplus T^{(0)}
               \oplus T^{(+1)}\oplus T^{(+2)},\quad
               (d_q)=(1,4,6,4,1).}
 \label{eq:blocks}
\end{equation}
Operators of the same order may mix. Hermiticity preservation,
$\Lambda_t(X^\dagger)=\Lambda_t(X)^\dagger$, relates the opposite orders,
since $E_{ab}^\dagger=E_{ba}$ has order $-q$. Using
Eq.~\eqref{eq:kraus}, $(K_{r\nu}E_{ab}K_{r\nu}^\dagger)_{cd}
=(K_{r\nu})_{ca}(K_{r\nu})_{db}^*$, so in the ordered basis $\mathcal B_q$
the exact entries are
\begin{equation}
 T^{(q)}_{cd,ab}(t)
 =\sum_{r,\nu}(K_{r\nu})_{ca}(K_{r\nu})_{db}^*,
 \qquad E_{ab},E_{cd}\in\mathcal B_q.
 \label{eq:blockentries}
\end{equation}
Equation~\eqref{eq:blockentries} is a formal construction, not a closed
evaluation of these functions for general interacting XXZ dynamics.

\subsection{Choi blocks}

By Eq.~\eqref{eq:choigeneral},
$J_t=\frac14\sum_{a,b}\sum_{c,d}T_{cd,ab}E_{cd}\otimes E_{ab}$, with
nonzero coefficients only for $m_c-m_d=m_a-m_b$. Since $M_S^{\mathsf T}=M_S$
is diagonal, $[M_S\otimes I_R-I_{AB}\otimes M_S^{\mathsf T},E_{cd}\otimes E_{ab}]
=[(m_c-m_d)-(m_a-m_b)]E_{cd}\otimes E_{ab}$, which vanishes on every
nonzero term. Hence the normalized Choi state satisfies
\begin{equation}
 [J_t,M_S\otimes I_R-I_{AB}\otimes M_S^{\mathsf T}]=0.
\end{equation}
Group output--reference kets $\ket c\ket a$ by $m_c-m_a$; the term
$E_{cd}\otimes E_{ab}=\ket{ca}\bra{db}$ connects $m_c-m_a$ to $m_d-m_b$,
which are equal. The numbers of pairs $(c,a)$ with $m_c-m_a=k$ are
$1,4,6,4,1$ for $k=-2,\ldots,2$, so the Choi blocks have the same sizes.
Since $J_t$ is Hermitian, the blocks contain at most
$1+16+36+16+1=70$ real parameters before trace preservation and positivity
are imposed, rather than $256$. Choi blocks and superoperator blocks are
different rearrangements of channel coefficients and should not be
identified as matrices or assigned the same spectra.

\subsection{Kraus operators with definite charge transfer}

Under Eq.~\eqref{eq:bulksym}, diagonalize $\omega_E$ within its charge
sectors: $M_E\ket{\phi_\nu}=\mu_\nu\ket{\phi_\nu}$. Choose
$M_E\ket r=\mu_r\ket r$. From $[U_t,M]=0$,
\begin{align}
 M_SK_{r\nu}&=\sqrt{p_\nu}\,{}_E\bra r(M-M_E)U_t\ket{\phi_\nu}_E
 =\sqrt{p_\nu}\,{}_E\bra rU_tM\ket{\phi_\nu}_E-\mu_rK_{r\nu}\notag\\
 &=\sqrt{p_\nu}\,{}_E\bra rU_t(M_S+\mu_\nu)\ket{\phi_\nu}_E-\mu_rK_{r\nu}
 =K_{r\nu}M_S+(\mu_\nu-\mu_r)K_{r\nu}.
\end{align}
Hence
\begin{equation}
 [M_S,K_{r\nu}]=(\mu_\nu-\mu_r)K_{r\nu},\qquad
 (K_{r\nu})_{ca}\ne0\ \Longrightarrow\
 m_c-m_a=\mu_\nu-\mu_r .
\end{equation}
Since $m_c-m_a\in\{-2,\ldots,2\}$, only transfers between $-2$ and $+2$
can act on the two probes; Kraus operators with other transfers vanish.
Under $V_\theta$, each $K_{r\nu}$ acquires the scalar phase
$e^{-i\theta(\mu_\nu-\mu_r)}$, which cancels in $K_{r\nu}XK_{r\nu}^\dagger$;
this reproduces Eq.~\eqref{eq:covariance}. The transfer label belongs to
a Kraus operator; coherence order labels an input or output matrix unit.
Their roles in the Choi and superoperator reductions are related but
distinct.

\subsection{Reduced states at fixed total magnetization}
\label{sec:xstate}

If the \emph{global} initial state is supported in one total-$M$ sector,
the reduced probe state has the exact form
\begin{equation}
 \rho_{AB}(t)=\begin{pmatrix}
 p_{00}&0&0&0\\0&p_{01}&c&0\\0&c^*&p_{10}&0\\0&0&0&p_{11}
 \end{pmatrix}.
 \label{eq:xstate}
\end{equation}
For a pure global ket $\ket{\Psi(t)}$ with $M\ket\Psi=m\ket\Psi$ (preserved
by $[U_t,M]=0$), resolve its four probe components,
$\ket{\Psi(t)}=\sum_a\ket a\otimes\ket{\psi_a(t)}$. Since $M_S\ket a=m_a\ket a$,
each bulk ket satisfies $M_E\ket{\psi_a}=(m-m_a)\ket{\psi_a}$; the charges are
$m-1,m,m,m+1$. The reduced state has entries
$\rho_{ab}=\langle\psi_b|\psi_a\rangle$, and bulk kets of different charge are
orthogonal. This eliminates every coherence except $01\leftrightarrow10$,
with $c=\langle\psi_{10}|\psi_{01}\rangle$. Linearity extends the conclusion
to mixtures within the sector.
The populations and $c$ remain time dependent. For this $X$ state with
$\rho_{00,11}=0$, the Wootters formula reduces to
\begin{equation}
 C=2\max\{0,|c|-\sqrt{p_{00}p_{11}}\}.
\end{equation}

This restriction does \emph{not} apply to the texture protocol merely
because the bulk has fixed magnetization. The input $\ket{++}$ contains
probe charges $+1,0,-1$ with probabilities $1/4,1/2,1/4$, and hence
occupies three total-charge sectors for a fixed-charge bulk. Their
coherences cannot be discarded. In particular, $f_1$ itself is not
of the form~\eqref{eq:xstate}.

Reflection and global spin flip give additional covariance only when
both the Hamiltonian and bulk preparation have the required symmetry.
Equal contacts or a symmetric Hamiltonian alone do not impose that
symmetry on an arbitrary reduced trajectory.

\subsection{Exchange texture and the limits of one input}
\label{sec:extexture}

Equations~\eqref{eq:textures}, \eqref{eq:measurement} and
\eqref{eq:response} hold for exchange coupling as well. However,
$\rho_f$ is generally not a Schur multiplier matrix and does not
determine the exchange channel. Under Eq.~\eqref{eq:bulksym},
$f_1=\frac14\sum_{a,b}E_{ab}$ has a nonzero component in every
$\mathcal V_q$, but one output determines only each block's action on
that particular component. It cannot determine a general $6\times6$ or
$4\times4$ block.

There is one exact scalar exception. Since $\mathcal V_{+2}$ is
one-dimensional, write $\Lambda_t(E_{00,11})=\xi_{+2}(t)E_{00,11}$. By the
block selection rule, the $E_{00,11}$ component of
$\rho_f=\frac14\sum_{a,b}\Lambda_t(E_{ab})$ comes only from
$\frac14\Lambda_t(E_{00,11})$. Then
\begin{equation}
 \xi_{+2}(t)=4[\rho_f(t)]_{00,11},\qquad
 \xi_{-2}(t)=\xi_{+2}(t)^*,
\end{equation}
the second relation by Hermiticity preservation.
This inference requires covariance; it is not guaranteed for a bulk
coherent across magnetization sectors or for transverse-field dynamics.
The equality $F_e=p_f$ is likewise not a general exchange-channel
identity, although $F_{\mathrm{avg}}=(4F_e+1)/5$ still holds.

\subsection{Applicability of the symmetry reduction}

The exchange block decomposition needs both
$[H,M_S+M_E]=0$ and $[\omega_E,M_E]=0$. Transverse bulk fields generally
violate total-magnetization conservation and couple different
total-magnetization sectors. Coherence between bulk magnetization
sectors generally removes channel phase covariance, even for a conserving
Hamiltonian. The Ising Schur channel and $\Gamma=4\rho_f$ remain valid in
both cases, because the probe $z$ components are still conserved.

\subsection{Summary of the exchange channel structure}

For exchange coupling, conservation of total magnetization and a
magnetization-invariant bulk preparation give five coherence-order
blocks. Symmetry fixes their dimensions and forbids mixing between
different orders, but leaves the entries within each block to be
calculated. The output from $\ket{++}$ determines the texture and the
one-dimensional highest-order response; it does not specify the full
exchange channel. A general closed expression for the interacting
exchange coefficients remains unavailable.

\appendix
\section{State diagnostics and their interpretation}\label{app:observables}

For $\rho=\rho_{AB}$, the state purity and entropy are
\begin{equation}
 P=\Tr\rho^2,\qquad S(\rho)=-\sum_k\lambda_k\ln\lambda_k,
 \qquad 0\ln0=0.
\end{equation}
Logarithms are natural. The probe mutual information is
$I(A:B)=S(\rho_A)+S(\rho_B)-S(\rho)$, with
$\rho_A=\Tr_B\rho$ and $\rho_B=\Tr_A\rho$. It includes classical
and quantum correlations.

Probe negativity and concurrence are
\begin{align}
 \mathcal N_{A|B}(\rho)&=\frac{\|\rho^{T_B}\|_1-1}{2},\\
 \widetilde\rho&=(\sigma_y\otimes\sigma_y)\rho^*
 (\sigma_y\otimes\sigma_y),\qquad
 C(\rho)=\max(0,r_1-r_2-r_3-r_4),
\end{align}
where the $r_i$ are the decreasing nonnegative square roots of the
eigenvalues of $\rho\widetilde\rho$ \cite{Wootters1998}.
For two qubits, $C>0$ and $\mathcal N>0$ are equivalent tests of
entanglement. They concern $A|B$, while $S(\rho_{AB})$ concerns the
pair's marginal mixedness. Only for a globally pure state does
$S(\rho_{AB})$ equal entanglement entropy across $AB|E$.

For the Ising texture input, Eq.~\eqref{eq:choipurity} with
$\sum_{a,b}|\Gamma_{ab}|^2=4+2\sum_{a<b}|\Gamma_{ab}|^2$ gives
\begin{equation}
 P_f(t)=\frac14+\frac18\sum_{a<b}|\Gamma_{ab}(t)|^2,
 \qquad \frac14\le P_f\le1.
\end{equation}
Its local probe coherences, obtained by tracing out the other probe, are
\begin{equation}
 (\rho_A)_{01}=\frac{\Gamma_{00,10}+\Gamma_{01,11}}4,
 \qquad (\rho_B)_{01}=\frac{\Gamma_{00,01}+\Gamma_{10,11}}4,
\end{equation}
and each local diagonal is $(1/2,1/2)$.
For a pure bulk, the global state is
$\frac12\sum_a\ket a\otimes U_a\ket{\phi_0}$; it is a product
across $AB|E$ iff its four conditional bulk kets are proportional.
For a mixed bulk, reduced entropy or purity alone does not quantify
entanglement across that cut.

\bibliographystyle{unsrtnat}
\bibliography{references}
\end{document}
